% GENERATED by blueprint/gen_content.py from the technical report; do not edit by hand.
This blueprint lays out the mathematics of the technical report \emph{Testing stemmata against the evidence: Technical report} (Metamorphoses of Civilization Collection, 2026, \href{https://doi.org/10.5281/zenodo.24405385}{doi:10.5281/zenodo.24405385}), Sections 3 to 9, and links every result to its proof in the Lean library \href{https://github.com/metamorphoses-collection/stemma-lean}{stemma-lean}. Statements and proofs are those of the report. Chapters carry the report's section numbers. The blueprint leaves out the report's remarks and examples, which the report numbers with its results, so within a chapter the numbers differ and each title gives the report's number. The dependency graph is that of the report's Appendix A: an arrow from one result to another means that the second result's proof uses the first. Every result and every definition is formalised and every proof is complete in Lean, with no \texttt{sorry} and only the standard axioms.
\setcounter{chapter}{2}

\chapter{Preliminaries}\label{sec:prelim}

\section{Splits and stemmata}

\begin{definition}[Splits and compatibility]
\label{def:split}
\lean{Stemma.Compat, Stemma.PairwiseCompat, Stemma.Splits.Nontriv}
\leanok
A \emph{split} of $Y\subseteq\X$ is an unordered bipartition $A|B$ with $A\cup B=Y$, $A\cap B=\emptyset$; it is \emph{full} if $Y=\X$, \emph{partial} otherwise, and \emph{non-trivial} if $|A|,|B|\ge2$. Two full splits $A|B$ and $C|D$ are \emph{compatible}, written $A|B\comp C|D$, if at least one of $A\cap C$, $A\cap D$, $B\cap C$, $B\cap D$ is empty, and \emph{in conflict} (they \emph{cross}) otherwise.
\end{definition}

\begin{definition}[Stemma]
\label{def:stemma}
\lean{Phylogenetics.SplitsEquivalence.IsXTree, Phylogenetics.SplitsEquivalence.splitSides, Stemma.Splits.IsSide, Stemma.Splits.IsSystem, Stemma.Splits.IsBinary}
\leanok
\uses{def:split}
An (unrooted) \emph{stemma} on $\X$ is a tree whose leaves are labelled bijectively by $\X$ and whose internal vertices have degree at least three. Removing an edge separates the leaves into a split; $\Sigma(T)$ is the set of non-trivial edge splits. $T$ is \emph{binary} if every internal vertex has degree three. Each split in $\Sigma(T)$ is a family of the stemma, the witnesses descending from one lost intermediary set against all the others, and a binary stemma is one in which every lost intermediary has exactly two lines of descent.
\end{definition}

\begin{proposition}[Buneman; Proposition 3.2 of the report]
\label{prop:buneman}
\lean{Stemma.Buneman.buneman}
\leanok
\uses{def:split, def:stemma}
A set $S$ of non-trivial full splits equals $\Sigma(T)$ for some stemma $T$ if and only if its splits are pairwise compatible; $T$ is then unique up to isomorphism.
\end{proposition}
\begin{proof}
\leanok
Classical (Buneman 1971; Semple and Steel 2003). The report states it without proof; the Lean library proves it in full.
\end{proof}

\begin{lemma}[Two counting facts; Lemma 3.3 of the report]
\label{lem:facts}
\lean{Stemma.Splits.card_le}
\leanok
\uses{def:split, def:stemma}
\begin{enumerate}[(i)]
\item A pairwise compatible set of non-trivial full splits has at most $n-3$ elements.
\item If $T$ is binary, $|\Sigma(T)|=n-3$, and every non-trivial split compatible with all of $\Sigma(T)$ belongs to $\Sigma(T)$.
\end{enumerate}
\end{lemma}
\begin{proof}
\leanok
\uses{prop:buneman}
(i) By Proposition~\ref{prop:buneman} the set is $\Sigma(T')$ for a stemma $T'$. A tree with $n$ leaves and $k$ internal vertices, all of degree at least three, has $n+k-1$ edges and degree sum $2(n+k-1)\ge n+3k$, so $k\le n-2$, and it has $(n+k-1)-n=k-1\le n-3$ internal edges, which carry its non-trivial splits. (ii) Equality holds in the degree count exactly when every internal vertex has degree three, so $|\Sigma(T)|=n-3$. If a non-trivial $\sigma\notin\Sigma(T)$ were compatible with all of $\Sigma(T)$, then $\Sigma(T)\cup\{\sigma\}$ would be a compatible set of $n-2$ splits, contradicting (i).
\end{proof}

\section{Collations and the estimator}

\begin{definition}[Collation, support and the best-supported stemma]
\label{def:collation}
\lean{Stemma.weight, Stemma.IsOptimal, Stemma.Splits.wt, Stemma.Splits.IsUniqueSolution}
\leanok
\uses{def:split}
A \emph{collation} of $\X$ is a finite sequence of \emph{variation units} $u=1,\dots,L$. Unit $u$ has a set $O_u\subseteq\X$ of witnesses legible at $u$ and a map $r_u:O_u\to R_u$ to a finite set of readings. It is \emph{complete} if $O_u=\X$ for all $u$.

Each reading $\rho$ of $u$ \emph{proposes} the split $r_u^{-1}(\rho)\,|\,O_u\setminus r_u^{-1}(\rho)$ of $O_u$. Let $P_u$ be the set of non-trivial splits proposed by $u$, each counted once (a binary unit proposes the same split from both readings). The \emph{support} or weight of a split is $w(\sigma)=|\{u:\sigma\in P_u\}|$, and for a set $S$, $w(S)=\sum_{\sigma\in S}w(\sigma)$. The \emph{candidates} are the splits of positive weight; the \emph{conflict graph} has the candidates as vertices and an edge between candidates in conflict. For a complete collation, the \emph{best-supported stemma} is
\begin{equation}\label{eq:mwis}
\hat S=\arg\max\{w(S):S\text{ a pairwise compatible set of candidates}\},
\end{equation}
a maximum-weight independent set of the conflict graph, and by Proposition~\ref{prop:buneman} a stemma; for characters, an exact maximum-compatibility algorithm of this kind is Cherry's~(Cherry 2017). When the maximum is not unique, $\hat S$ denotes any maximiser; every statement below holds for every selection rule unless ties are discussed explicitly. With no candidates, $\hat S=\emptyset$, the unresolved star. A multistate unit is coded by its reading-versus-rest splits, each counted once per unit. For binary units this is classical maximum compatibility with repeated characters; for multistate units it is not the compatibility of whole multistate characters.
\end{definition}

\begin{lemma}[Exclusivity; Lemma 3.5 of the report]
\label{lem:exclusive}
\lean{Stemma.exclusivity}
\leanok
\uses{def:collation}
Two splits proposed by the same unit are compatible. Equivalently, splits in conflict are never proposed by the same unit.
\end{lemma}
\begin{proof}
\leanok
For distinct readings $\rho,\rho'$ of $u$, the sides $r_u^{-1}(\rho)$ and $r_u^{-1}(\rho')$ are disjoint, so one of the four intersections is empty.
\end{proof}

\section{The rooted variant}

\begin{definition}[Clusters and the rooted variant]
\label{def:rooted}
\lean{Stemma.Splits.Laminar, Stemma.Rooted.Lam, Stemma.Rooted.IsCluster, Stemma.Rooted.IsBinR}
\leanok
\uses{def:collation}
When a reference reading $\rho_u^\ast$ is fixed in each unit (an editor's lemma, a majority reading), a unit proposes \emph{clusters}, the sets $r_u^{-1}(\rho)$ with $\rho\neq\rho^\ast_u$ and $|r_u^{-1}(\rho)|\ge2$: the witnesses sharing a departure from the reference. Two clusters are \emph{compatible} if they are disjoint or nested (\emph{laminar}). This is the common-error method in its classical form: only agreement in a departure from the reference groups witnesses, and agreement in the reference reading proves nothing. A rooted stemma is determined by its clusters of size at least two~(Semple and Steel 2003), and~\eqref{eq:mwis} applies verbatim to clusters. Clusters from one unit are disjoint, hence compatible: Lemma~\ref{lem:exclusive} holds for clusters.
\end{definition}

\chapter{Complexity}\label{sec:complexity}

\begin{theorem}[Universality; Theorem 4.1 of the report]
\label{thm:universal}
\lean{Stemma.Universal.universality}
\leanok
\uses{def:split}
Every finite simple graph $H=(V,E)$ with $|V|\ge2$ is the conflict graph of a family of distinct non-trivial full splits on $4|E|+2|V|+2$ witnesses.
\end{theorem}
\begin{proof}
\leanok
For each vertex $v$ create two \emph{private} witnesses; for each edge $e=\{u,v\}$, with its endpoints in a fixed order, create four witnesses $q_e^{ab}$, $a,b\in\{0,1\}$, where $a$ is the coordinate of the first endpoint and $b$ that of the second; add two auxiliary witnesses $z_1,z_2$. Let side~1 of $\sigma_v$ consist of the private witnesses of $v$ and every $q_e^{ab}$ with $e\ni v$ whose coordinate for $v$ is~$1$; side~0 is everything else, including $z_1,z_2$. Side~1 has at least the two private witnesses and side~0 at least $z_1,z_2$.

The splits are distinct. Side~1 of $\sigma_u$ contains no $z_i$, so it can equal only side~1 of $\sigma_v$, which contains the private witnesses of $v$ and not those of $u$.

If $e=\{u,v\}\in E$, the four witnesses $q_e^{11},q_e^{10},q_e^{01},q_e^{00}$ lie respectively on side~1 of both, side~1 of the first endpoint only, side~1 of the second only, and side~0 of both; all four intersections are non-empty, so $\sigma_u$ and $\sigma_v$ conflict. If $\{u,v\}\notin E$, side~1 of $\sigma_u$ and side~1 of $\sigma_v$ share no private witness, and no edge witness, because an edge witness lies on side~1 of $\sigma_x$ only if $x$ is an endpoint of its edge and no edge has both $u$ and $v$ as endpoints. The sides are disjoint and the splits compatible.
\end{proof}

\begin{corollary}[Corollary 4.2 of the report]
\label{cor:nphard}
\lean{Stemma.NPHard.CompatibleSplits_npComplete, Stemma.NPHard.WeightedSplits_npComplete}
\leanok
\uses{def:collation}
MWIS restricted to split-conflict graphs is NP-hard when $n$ is part of the input. Unless P$=$NP, no class of graphs on which MWIS is solvable in polynomial time contains every split-conflict graph.
\end{corollary}
\begin{proof}
\leanok
\uses{thm:universal}
The construction is polynomial in $|V|+|E|$ and preserves vertex weights, so it reduces MWIS on general graphs, which is NP-hard~(Karp 1972), to MWIS on split-conflict graphs.
\end{proof}

\chapter{Fragmentary witnesses}\label{sec:partial}

\begin{definition}[Partial splits and overlap-compatibility]
\label{def:partial}
\lean{Stemma.PSplit, Stemma.overlapCompat}
\leanok
\uses{def:split}
At such places a unit proposes \emph{partial} splits of $O_u$. Two partial splits $A|B$ of $Y$ and $C|D$ of $Z$ are \emph{overlap-compatible} if their restrictions to $Y\cap Z$ are compatible, with the convention that a restriction with an empty side is compatible with everything.
\end{definition}

\begin{definition}[Joint completability]
\label{def:joint}
\lean{Stemma.Extends, Stemma.JointPair, Stemma.Cutting.JointlyCompletable}
\leanok
\uses{def:partial}
A family of partial splits is \emph{jointly completable} if each can be extended to a full split of $\X$ so that the extensions are pairwise compatible.
\end{definition}

\begin{proposition}[Pairs are exact, triples are not; Proposition 5.2 of the report]
\label{prop:triple}
\lean{Stemma.pair_exact, Stemma.triple_not_joint}
\leanok
\uses{def:partial, def:joint}
\begin{enumerate}[(a)]
\item Two overlap-compatible partial splits are jointly completable.
\item Three partial splits can be pairwise overlap-compatible without being jointly completable: on $\X=\{1,\dots,5\}$, $p=12|34$ (on $1234$), $q=12|35$ (on $1235$), $r=14|25$ (on $1245$).
\item Five witnesses is the minimum: on four witnesses no such triple exists.
\end{enumerate}
\end{proposition}
\begin{proof}
\leanok
\uses{prop:buneman}
(a) Let $A|B$ on $Y$ and $C|D$ on $Z$ be compatible on $Y\cap Z$; up to renaming sides, $A\cap C\cap Y\cap Z=\emptyset$. Since $A\subseteq Y$ and $C\subseteq Z$, $A\cap C=\emptyset$. Put every unobserved witness on the side opposite $A$, respectively $C$: the full splits $A|\X\setminus A$ and $C|\X\setminus C$ are compatible because $A\cap C=\emptyset$.

(b) On the overlaps, $p$ and $q$ both restrict to $12|3$; $p$ and $r$ restrict to $12|4$ and $14|2$; $q$ and $r$ restrict to $12|5$ and $1|25$. In each pair one intersection is empty, so the three are pairwise overlap-compatible. A completion places witness $5$ in $p$, $4$ in $q$ and $3$ in $r$, so there are $2^3$ completions. Writing each completion as the side that contains witness $1$, $p\in\{12,125\}$, $q\in\{12,124\}$ and $r\in\{14,134\}$. If $r=14$, it conflicts with $p=12$ (the four intersections are $1$, $4$, $2$, $35$) and with $p=125$ ($1$, $4$, $25$, $3$). If $r=134$, it conflicts with $q=12$ ($1$, $34$, $2$, $5$) and with $q=124$ ($14$, $3$, $2$, $5$). Every completion therefore contains a conflicting pair (Appendix~B of the report).

(c) A non-trivial partial split has at least four observed witnesses, so on four witnesses every non-trivial split is full, and pairwise compatibility is joint compatibility by Proposition~\ref{prop:buneman}.
\end{proof}

\begin{theorem}[Exact selection by cutting planes; Theorem 5.3 of the report]
\label{thm:cutting}
\lean{Stemma.CuttingPlanes.cutting_correct}
\leanok
\uses{def:joint, def:collation}
Write~\eqref{eq:mwis} on the overlap-conflict graph as an integer programme, with a variable $y_\sigma\in\{0,1\}$ for each candidate and a constraint $y_\sigma+y_{\sigma'}\le1$ for each edge, and solve it together with the cuts already added. If the optimum $S$ is jointly completable, stop. Otherwise find a non-completable $K\subseteq S$, add the cut $\sum_{\sigma\in K}y_\sigma\le|K|-1$, and repeat. The procedure terminates, and its first jointly completable optimum is a maximum-weight jointly completable family.
\end{theorem}
\begin{proof}
\leanok
A jointly completable family is pairwise overlap-compatible, since compatible full splits restrict to compatible splits on any subset; so the first master problem is a relaxation. Every subfamily of a jointly completable family is jointly completable (keep the same extensions), so a cut excludes only families containing $K$, none of which is completable; every master remains a relaxation and its optimum an upper bound. Each cut removes the current optimum and there are finitely many subsets, so the loop terminates; a completable master optimum is feasible and attains the bound.
\end{proof}

\chapter{Deterministic certificates of uniqueness}\label{sec:certificates}

\begin{lemma}[Displacement matching; Lemma 6.1 of the report]
\label{lem:hall}
\lean{Stemma.Splits.displacement_matching}
\leanok
\uses{def:stemma}
Let $T$ be binary, $S\neq\T$ a pairwise compatible set of non-trivial splits, $A=\T\setminus S$ and $B=S\setminus\T$. Then $A\neq\emptyset$, $|B|\le|A|$, and there is an injection $\varphi:B\to A$ such that $\tau$ and $\varphi(\tau)$ are in conflict for every $\tau\in B$.
\end{lemma}
\begin{proof}
\leanok
\uses{lem:facts}
A split $\tau\in B$ is compatible with $S\cap\T$, so every true split in conflict with $\tau$ lies in $A$; by Lemma~\ref{lem:facts}(ii) there is at least one. For $B'\subseteq B$ let $N(B')\subseteq A$ be the set of true splits in conflict with some member of $B'$. The set $(\T\setminus N(B'))\cup B'$ is pairwise compatible: true splits are pairwise compatible, members of $B'\subseteq S$ are, and each member of $B'$ is compatible with every true split outside $N(B')$ by definition of $N$. By Lemma~\ref{lem:facts}(i) it has at most $n-3=|\T|$ elements; since $B'\cap\T=\emptyset$ this reads $|\T|-|N(B')|+|B'|\le|\T|$, that is $|B'|\le|N(B')|$. Hall's theorem~(Hall 1935) gives a matching of $B$ into $A$ along conflicts. If $A=\emptyset$, then $S\supsetneq\T$, so $S$ contains a non-trivial split outside $\T$ compatible with all of $\T$, contradicting Lemma~\ref{lem:facts}(ii).
\end{proof}

\begin{theorem}[Local separation; Theorem 6.2 of the report]
\label{thm:local}
\lean{Stemma.Splits.local_separation}
\leanok
\uses{def:stemma, def:collation}
Let $T$ be binary, and let $w\ge0$ be defined on all non-trivial splits (zero off the candidates). If $w(\tau)<w(\sigma)$ whenever a false split $\tau$ and a true split $\sigma$ are in conflict, then $\T$ is the unique solution of~\eqref{eq:mwis}.
\end{theorem}
\begin{proof}
\leanok
\uses{lem:hall}
Every true split has positive weight. Indeed, an internal edge of a binary tree separates four non-empty subtrees $A_1,A_2$ on one side and $B_1,B_2$ on the other, and the non-trivial split $A_1\cup B_1|A_2\cup B_2$ is in conflict with it, hence false; the hypothesis gives $w(\sigma)>w(\tau)\ge0$. So $\T$ is a set of candidates. Let $S\neq\T$ be compatible, with $A$, $B$, $\varphi$ as in Lemma~\ref{lem:hall}. Then
\[
w(\T)-w(S)=w(A)-w(B)=\sum_{\tau\in B}\bigl(w(\varphi(\tau))-w(\tau)\bigr)+w\bigl(A\setminus\varphi(B)\bigr).
\]
If $B\neq\emptyset$, every bracket is positive and the last term non-negative. If $B=\emptyset$, the expression is $w(A)>0$ because $A\neq\emptyset$ and true splits have positive weight.
\end{proof}

\begin{theorem}[Global separation; Theorem 6.3 of the report]
\label{thm:global}
\lean{Stemma.Splits.global_separation}
\leanok
\uses{def:stemma, def:collation}
Let $T$ be binary. If $\min_{\sigma\in\T}w(\sigma)>\max_{\tau\notin\T}w(\tau)$, then $\T$ is the unique solution of~\eqref{eq:mwis}.
\end{theorem}
\begin{proof}
\leanok
\uses{thm:local}
The hypothesis implies that of Theorem~\ref{thm:local}.
\end{proof}

\begin{proposition}[Threshold for non-binary stemmata; Proposition 6.9 of the report]
\label{prop:threshold}
\lean{Stemma.Splits.threshold}
\leanok
\uses{def:stemma, def:collation}
For any stemma $T$ with at least one non-trivial split, if $\max_{\tau\notin\T}w(\tau)\le t<\min_{\sigma\in\T}w(\sigma)$, the solution of~\eqref{eq:mwis} restricted to splits of weight greater than $t$ is $\T$.
\end{proposition}
\begin{proof}
\leanok
Only true splits survive the threshold; they are pairwise compatible and have positive weight, so all are selected.
\end{proof}

\begin{proposition}[Rooted analogue; Proposition 6.10 of the report]
\label{prop:rooted}
\lean{Stemma.Rooted.isBinR_iff_outgroup, Stemma.RootedConsistency.consistency_rooted}
\leanok
\uses{def:rooted}
Lemma~\ref{lem:hall}, Theorems~\ref{thm:local} and~\ref{thm:global}, Lemma~\ref{lem:gamma} and Theorem~\ref{thm:consistency} hold for clusters, with $T$ a binary rooted stemma, $\T$ its \emph{proper} clusters (of size between $2$ and $n-1$), conflict meaning not laminar, compatible sets replaced by laminar families of proper clusters, the margin $\gamma$ of Definition~\ref{def:gamma} taken over laminar families $S\neq\T$, and $n-3$ replaced by $n-2$ in the bound of Theorem~\ref{thm:consistency}(a), in the complete-observation setting.
\end{proposition}
\begin{proof}
\leanok
\uses{lem:facts, lem:hall, thm:local, thm:global, lem:gamma, thm:consistency}
The proofs use only Lemma~\ref{lem:facts}, with $n-3$ replaced by $|\T|$, and, for the positivity of true supports in Theorem~\ref{thm:local}, a false cluster crossing each true cluster $C$: with $C_1$ a child of $C$ and $D$ the sibling of $C$, the cluster $C_1\cup D$ crosses $C$. For proper clusters its two facts read: (i) a laminar family of distinct subsets of $\X$ of sizes between $2$ and $n-1$ has at most $n-2$ members; (ii) the proper clusters of a binary rooted stemma are $n-2$ and form a maximal such family. For (i), add all singletons and $\X$; in the enlarged laminar family each set of size at least two is the disjoint union of its maximal proper members (an element not covered by a non-singleton proper member is itself a singleton member), of which there are at least two, so the Hasse diagram is a tree with $n$ leaves in which every internal vertex has at least two children; it has at most $n-1$ internal vertices, one of which is $\X$. For (ii), a binary rooted tree on $n$ leaves has $n-1$ internal vertices, one cluster each, and the root's cluster is $\X$. Let $C$ with $2\le|C|\le n-1$ be laminar with every cluster, and let $V$ be the smallest cluster containing $C$, with children $V_1,V_2$. If $C\subseteq V_i$, either $|V_i|\ge2$, contradicting minimality, or $|V_i|=1$, contradicting $|C|\ge2$. So for each $i$, $V_i\subseteq C$ or $V_i\cap C=\emptyset$, and $C$ is $V_1$, $V_2$ or $V$; it is a proper cluster since $C\ne\X$. Lemma~\ref{lem:gamma} and Theorem~\ref{thm:consistency} use, beyond the definitions, only that $\T\setminus S\ne\emptyset$ and $|S\setminus\T|\le|\T\setminus S|$ for every compatible $S\neq\T$. For a laminar family $S\neq\T$ of proper clusters, (i) and (ii) give $|S|\le n-2=|\T|$, hence $|S\setminus\T|=|S|-|S\cap\T|\le|\T|-|S\cap\T|=|\T\setminus S|$; and $\T\setminus S=\emptyset$ would give $S\supsetneq\T$, against the maximality in (ii). With $m$, $w$, $\F$ and $K$ defined on clusters, the two proofs then apply verbatim, the union bound running over the $n-2+K$ clusters that can occur.
\end{proof}

\chapter{Consistency and sample size}\label{sec:consistency}

\section{A pairwise bound}

\begin{theorem}[Pairwise bound; Theorem 7.1 of the report]
\label{thm:pairwise}
\lean{Stemma.Chernoff.pairwise_bound}
\leanok
\uses{def:gamma}
Let $T$ be binary, $m_\sigma>0$ for every true $\sigma$, and $m_\tau<m_\sigma$ whenever $\tau\in\F$ and $\sigma\in\T$ are in conflict. Then
\begin{align*}
\Pr(\hat S_L\neq\T)\le{}&\sum_{\substack{(\tau,\sigma)\in\F\times\T\\\text{in conflict}}}\bigl(1-(\sqrt{m_\sigma}-\sqrt{m_\tau})^2\bigr)^L\\
&+\sum_{\sigma\in\T}(1-m_\sigma)^L .
\end{align*}
\end{theorem}
\begin{proof}
\leanok
\uses{lem:exclusive, thm:local}
By Theorem~\ref{thm:local}, applied with $w$ extended by zero, failure implies that some conflicting pair has $w(\tau)\ge w(\sigma)$, which for $\tau\notin\F$ means $w(\sigma)=0$; ties are counted as failures, whatever the selection rule. The event $w(\sigma)=0$ has probability $(1-m_\sigma)^L$, which gives the second sum. For a pair with $\tau\in\F$, the first sum, by Lemma~\ref{lem:exclusive} the difference $w(\sigma)-w(\tau)$ is a sum of $L$ independent variables taking the values $1,-1,0$ with probabilities $m_\sigma,m_\tau,1-m_\sigma-m_\tau$. For $t>0$, Markov's inequality applied to $e^{-t(w(\sigma)-w(\tau))}$ gives $\Pr(w(\sigma)-w(\tau)\le0)\le\bigl(m_\sigma e^{-t}+m_\tau e^{t}+1-m_\sigma-m_\tau\bigr)^L$~(Chernoff 1952); at $e^{t}=\sqrt{m_\sigma/m_\tau}$ the bracket equals $1-(\sqrt{m_\sigma}-\sqrt{m_\tau})^2$. A union bound concludes.
\end{proof}

\section{The margin and the consistency theorem}

\begin{definition}[Margin; Definition 7.2 of the report]
\label{def:gamma}
\lean{Stemma.Margin.competitors, Stemma.Margin.ratio, Stemma.Margin.margin}
\leanok
\uses{def:collation, def:stemma}
Let the units be independent and identically distributed under some copying process, with $n$ fixed; nothing else is assumed. Let $m_\sigma=\Pr(\sigma\in P_u)$, $m(S)=\sum_{\sigma\in S}m_\sigma$, $\F=\{\tau\notin\T:m_\tau>0\}$, $K=|\F|$, and let $\hat S_L$ be a solution of~\eqref{eq:mwis} on $L$ units. Each $w(\sigma)$ is binomial with parameters $L$ and $m_\sigma$.

For binary $T$, the \emph{margin} is
\[
\gamma=\min_{S\neq\T\text{ compatible}}\frac{m(\T\setminus S)-m(S\setminus\T)}{|\T\setminus S|}.
\]
\end{definition}

\begin{lemma}[Computing $\gamma$; Lemma 7.3 of the report]
\label{lem:gamma}
\lean{Stemma.Margin.computing_gamma}
\leanok
\uses{def:gamma}
For $g\in\mathbb{R}$ let $m^{(g)}_\sigma=m_\sigma-g$ for $\sigma\in\T$ and $m^{(g)}_\tau=m_\tau$ otherwise. Then $\T$ is a maximum of $m^{(g)}$ over compatible sets if and only if $g\le\gamma$. Hence, given the expected supports, $\gamma$ is approximated to any tolerance by bisection on $g$ in the bracket $[-1,1]$ (since $|B|\le|A|$ and $0\le m_\sigma\le1$), with one MWIS per step, each step comparing optimum values rather than relying on the solver to return $\T$ at a tie. A value computed from data is an estimate of $\gamma$, not $\gamma$.
\end{lemma}
\begin{proof}
\leanok
\uses{lem:hall}
For compatible $S\neq\T$, $m^{(g)}(\T)-m^{(g)}(S)=m(A)-m(B)-g|A|$ with $A=\T\setminus S$, $B=S\setminus\T$, since $g$ is subtracted once for each true split in $\T$ and in $S$, and those in both cancel. This is non-negative for all $S$ if and only if $g\le\min_S(m(A)-m(B))/|A|=\gamma$. The set of such $g$ is a half-line, which makes bisection valid.
\end{proof}

\begin{theorem}[Consistency; Theorem 7.4 of the report]
\label{thm:consistency}
\lean{Stemma.Consistency.consistency, Stemma.Margin.inconsistency, Stemma.Bernstein.consistency_bernstein}
\leanok
\uses{def:gamma}
Let $T$ be binary.
\begin{enumerate}[(a)]
\item If $\gamma>0$, then
\[
\Pr(\hat S_L\neq\T)\le\sum_{\sigma\in\T\cup\F}2\exp\Bigl(-\frac{L\gamma^2}{8m_\sigma+4\gamma/3}\Bigr),
\qquad
\Pr(\hat S_L\neq\T)\le2(n-3+K)\,e^{-L\gamma^2/2}.
\]
\item If some compatible $S$ has $m(S)>m(\T)$, then $\Pr(\hat S_L=\T)\to0$ as $L\to\infty$.
\end{enumerate}
When $\T$ is a maximum of $m$ but not the unique one, the outcome depends on the selection rule and on the joint law of the supports, and neither consistency nor inconsistency follows (Remark~7.5 of the report).
\end{theorem}
\begin{proof}
\leanok
\uses{lem:hall}
(a) Suppose $|w(\sigma)/L-m_\sigma|<\gamma/2$ for every $\sigma\in\T\cup\F$; every other split has $w=m=0$ almost surely. For compatible $S\neq\T$, with $A,B$ as in Lemma~\ref{lem:hall} and $|B|\le|A|$,
\[
\frac{w(A)-w(B)}{L}>m(A)-m(B)-\frac\gamma2\bigl(|A|+|B|\bigr)\ge m(A)-m(B)-\gamma|A|\ge0,
\]
the last step by the definition of $\gamma$. Taking $S=\T\setminus\{\sigma\}$ in the definition gives $\gamma\le m_\sigma$ for every true $\sigma$, so on this event every true split has $w(\sigma)/L>m_\sigma-\gamma/2>0$ and $\T$ is a set of candidates. So $\T$ is the unique solution. By Bernstein's inequality~(Boucheron et al. 2013) for a binomial, with variance $Lm_\sigma(1-m_\sigma)\le Lm_\sigma$ and centred increments bounded by $1$, $\Pr(|w(\sigma)/L-m_\sigma|\ge t)\le2\exp\{-Lt^2/(2(m_\sigma+t/3))\}$; at $t=\gamma/2$ this is the stated term. Hoeffding's inequality~(Hoeffding 1963) gives $2e^{-2Lt^2}=2e^{-L\gamma^2/2}$ for each of the $n-3+K$ splits. A union bound concludes.

(b) $w(S)-w(\T)=\sum_u\bigl(|P_u\cap S|-|P_u\cap\T|\bigr)$ is a sum of $L$ independent bounded terms with positive mean $m(S)-m(\T)$, so $\Pr(w(S)>w(\T))\to1$ by the weak law of large numbers, and on that event $\T$ is not optimal.
\end{proof}

\begin{corollary}[Sample size; Corollary 7.6 of the report]
\label{cor:samplesize}
\lean{Stemma.Margin.sample_size}
\leanok
\uses{def:gamma}
If $\gamma>0$, $0<\varepsilon<1$ and $m_{\max}=\max_{\sigma\in\T\cup\F}m_\sigma$, then $\Pr(\hat S_L\neq\T)\le\varepsilon$ as soon as
\[
L\ge\frac{8m_{\max}+4\gamma/3}{\gamma^2}\,\log\frac{2(n-3+K)}{\varepsilon}.
\]
\end{corollary}
\begin{proof}
\leanok
\uses{thm:consistency}
The report gives no separate proof: it follows from Theorem~\ref{thm:consistency}. The Lean library proves it in full.
\end{proof}

\begin{proposition}[Dependent blocks; Proposition 7.8 of the report]
\label{prop:blocks}
\lean{Stemma.Blocks.dependent_blocks}
\leanok
\uses{def:gamma}
Let $T$ be binary. Suppose the units form $J$ independent, identically distributed blocks of $k$ consecutive units each ($L=Jk$), with arbitrary dependence inside a block, let $m_\sigma=\E[w_j(\sigma)]/k$, where $w_j$ is the support within block $j$, and define $\gamma$ from this $m$. If $\gamma>0$, then $\Pr(\hat S_L\neq\T)\le2(n-3+K)e^{-J\gamma^2/2}$.
\end{proposition}
\begin{proof}
\leanok
\uses{thm:consistency}
The proof of (a) uses only the deviations $|w(\sigma)/L-m_\sigma|$. Here $w(\sigma)/L=J^{-1}\sum_j w_j(\sigma)/k$ is an average of $J$ independent variables in $[0,1]$ with mean $m_\sigma$, and Hoeffding's inequality gives $2e^{-2J(\gamma/2)^2}$.
\end{proof}

\section{Certifying the margin from the data}\label{sec:gammaconf}

\begin{theorem}[Confidence bound for the margin; Theorem 7.9 of the report]
\label{thm:gammaconf}
\lean{Stemma.GammaConf.gamma_confidence}
\leanok
\uses{def:gamma}
Let units be i.i.d.; a unit may propose several splits, each counted once. For each of the $N=2^{n-1}-n-1$ non-trivial splits $s$, $w(s)$ is binomial with parameters $L$ and $m_s$; let $\ell_s$ and $u_s$ be its one-sided Clopper--Pearson lower and upper bounds at level $\zeta/(2N)$. For a binary stemma $T$ with split set $\T$ put
\[
\gamma_{\mathrm{lo}}(T)=\min_{S\neq\T\ \text{compatible}}\frac{\ell(\T\setminus S)-u(S\setminus\T)}{|\T\setminus S|}.
\]
Then, with probability at least $1-\zeta$, $\gamma(T)\ge\gamma_{\mathrm{lo}}(T)$ for every binary stemma $T$ simultaneously, in particular for $\hat S_L$ when it is binary. Consequently, the rule ``certify $\hat S_L$ when it is binary and $\gamma_{\mathrm{lo}}(\hat S_L)>0$'' issues a false certificate, one for a stemma that is not the unique maximum of expected support, with probability at most $\zeta$.
\end{theorem}
\begin{proof}
\leanok
\uses{lem:gamma}
Each of the $2N$ one-sided bounds fails with probability at most $\zeta/(2N)$, so all hold together with probability at least $1-\zeta$; call this event $\mathcal E$. On $\mathcal E$, for every binary $T$ and every compatible $S\neq\T$, where $\T\setminus S\neq\emptyset$ because $\T$ is a maximal compatible set, $m(\T\setminus S)-m(S\setminus\T)\ge\ell(\T\setminus S)-u(S\setminus\T)$, and taking the minimum gives $\gamma(T)\ge\gamma_{\mathrm{lo}}(T)$. The event $\mathcal E$ does not depend on $T$, so the inequality holds for a stemma chosen from the data. If $\gamma(\hat S_L)>0$, the split set of $\hat S_L$ is the unique maximum of $m$ (Lemma~\ref{lem:gamma}). The margin is defined for binary stemmata only: for a non-binary $T$ a compatible refinement $S\supsetneq\T$ has $\T\setminus S=\emptyset$, and a non-binary split set cannot be the unique maximum of non-negative weights.
\end{proof}

\begin{proposition}[Certifying the margin with dependent units; Proposition 7.11 of the report]
\label{prop:gammablocks}
\lean{Stemma.GammaBlocks.gamma_confidence_blocks}
\leanok
\uses{def:gamma}
Let the units form $J$ independent blocks of $k$ consecutive units ($L=Jk$), with arbitrary dependence inside a block and the same expected support in every block, and let $m_s=\E[w_j(s)]/k$, where $w_j(s)$ is the support of $s$ within block $j$. With $h=\sqrt{\log(2N/\zeta)/(2J)}$, put $\ell_s=w(s)/L-h$ and $u_s=w(s)/L+h$, and define $\gamma_{\mathrm{lo}}$ as in Theorem~\ref{thm:gammaconf}. Then the conclusions of Theorem~\ref{thm:gammaconf} hold.
\end{proposition}
\begin{proof}
\leanok
\uses{thm:gammaconf, prop:blocks}
$w(s)/L=J^{-1}\sum_j w_j(s)/k$ is an average of $J$ independent variables in $[0,1]$ with mean $m_s$, so by Hoeffding's inequality each one-sided bound fails with probability at most $e^{-2Jh^2}=\zeta/(2N)$. The rest of the proof of Theorem~\ref{thm:gammaconf} uses only that all $2N$ bounds hold together.
\end{proof}

\section{Consistency with fragmentary witnesses}\label{sec:fragconsistency}

\begin{theorem}[Fragmentary witnesses; Theorem 7.12 of the report]
\label{thm:fragmentary}
\lean{Stemma.FragmentaryUnits.fragmentary_consistency}
\leanok
\uses{def:gamma, def:partial}
Let $T$ be binary.
\begin{enumerate}[(a)]
\item Theorem~\ref{thm:consistency} holds for the completion estimator with $\tilde m$ in place of $m$.
\item For every full split, $\tilde m_\sigma\ge(1-4\delta)\,m_\sigma$ and $\tilde m_\tau\le m_\tau+2\sum_{\sigma'\neq\tau}m_{\sigma'}\,\delta^{d(\sigma',\tau)}$.
\item Consequently, if $m_\sigma-m_\tau\ge\varrho$ whenever a false $\tau$ and a true $\sigma$ are in conflict, the completion estimator is consistent for every $\delta<\varrho/(4+2\mu)$.
\end{enumerate}
\end{theorem}
\begin{proof}
\leanok
\uses{lem:hall, thm:local, thm:consistency}
(a) The proof of Theorem~\ref{thm:consistency} uses only that units are i.i.d., that each $\tilde w(\sigma)$ counts each unit at most once, so that $\tilde w(\sigma)/L$ is an average of i.i.d.\ indicators with mean $\tilde m_\sigma$, and Lemma~\ref{lem:hall}, which concerns full splits only. (The pairwise bound of Theorem~\ref{thm:pairwise} does not carry over: two conflicting full splits can share a restriction.)

(b) If a complete unit proposes $\sigma$ and $\sigma|_{O_u}$ is non-trivial, the unit proposes $\sigma|_{O_u}$. Fix two witnesses on each side of $\sigma$; if all four are legible, the restriction is non-trivial, and the probability that one of them is not is at most $4\delta$. By independence of $O_u$ from the readings, $\tilde m_\sigma\ge(1-4\delta)m_\sigma$. Conversely, a unit proposes $\tau|_{O_u}$ only if some complete proposal $\sigma'$ satisfies $\sigma'|_{O_u}=\tau|_{O_u}$. For $\sigma'=\tau$ this contributes at most $m_\tau$. For $\sigma'\neq\tau$, the restrictions agree only if every witness on which the two splits disagree, in one of the two orientations, is illegible, which has probability at most $\delta^{|A\triangle C|}+\delta^{|A\triangle D|}\le2\delta^{d(\sigma',\tau)}$; summing over $\sigma'$ gives the upper bound.

(c) Since $d\ge1$ for $\sigma'\ne\tau$ and $m_\sigma\le1$, part (b) gives $\tilde m_\sigma-\tilde m_\tau\ge m_\sigma-m_\tau-4\delta m_\sigma-2\delta\mu\ge\varrho-\delta(4+2\mu)>0$ for every conflicting pair. By Theorem~\ref{thm:local} applied to $\tilde m$, $\T$ is the unique maximum of $\tilde m$, so $\tilde\gamma>0$, and part (a) concludes.
\end{proof}

\begin{proposition}[Illegibility that depends on the text; Proposition 7.13 of the report]
\label{prop:textillegible}
\lean{Stemma.TextIllegible.margin_pos}
\leanok
\uses{def:gamma, def:partial}
Suppose the pairs (readings, legible set) of the units are independent and identically distributed, and that, given the splits a unit proposes in the complete data, each witness is illegible with probability at most $\delta$ (in particular, this holds if it holds given the readings); legibility may otherwise depend on the readings, on the text and on the other witnesses of the unit. Then Theorem~\ref{thm:fragmentary}(a) holds, the bounds of Theorem~\ref{thm:fragmentary}(b) hold with $\delta^{d(\sigma',\tau)}$ replaced by $\delta$, and the sufficient conditions of Theorem~\ref{thm:fragmentary}(c) and Theorem~\ref{thm:cuttingconsistency}(c) hold unchanged.
\end{proposition}
\begin{proof}
\leanok
\uses{thm:fragmentary}
Part (a) of Theorem~\ref{thm:fragmentary} uses only that units are i.i.d. In part (b), the lower bound needs four fixed witnesses to be legible; given the complete proposals, the probability that one of them is not is at most $4\delta$ by a union bound, whatever their joint law, so $\tilde m_\sigma\ge(1-4\delta)m_\sigma$ after averaging. For the upper bound, $\sigma'|_{O_u}=\tau|_{O_u}$ with $\sigma'\ne\tau$ requires, in one of the two orientations, a non-empty mismatch set to be illegible, hence one fixed witness of it; given the complete proposals this has probability at most $\delta$ for each orientation, so $\tilde m_\tau\le m_\tau+2\delta\mu$. Part (c) used only these two bounds (with $\delta^d\le\delta$). In Theorem~\ref{thm:cuttingconsistency}(c) every illegibility event is likewise one in which a fixed non-empty set, hence one fixed witness, is illegible, or a union of four such events, so its bound holds as stated.
\end{proof}

\begin{theorem}[Consistency of the cutting-plane estimator; Theorem 7.14 of the report]
\label{thm:cuttingconsistency}
\lean{Stemma.Cutting.cutting_consistent, Stemma.CuttingIllegible.cutting_illegible_margin}
\leanok
\uses{def:gamma, def:joint}
\begin{enumerate}[(a)]
\item A family of partial splits is jointly completable if and only if some binary stemma displays every member. Hence the cutting-plane estimator selects the observed partial splits displayed by a binary stemma $\hat T$ that maximises the total weight $W(T')$ of the observed partial splits it displays.
\item With i.i.d.\ observed units, $\hat T=T$ with probability tending to one if $T$ is the unique maximum of $M$ over binary stemmata, and with probability tending to zero if some $T'$ has $M(T')>M(T)$; in the first case the selected family is eventually displayed by $T$ alone. At a tie between maximisers nothing is claimed.
\item Under the illegibility model of Theorem~\ref{thm:fragmentary}, for every binary $T'\neq T$,
\[
M(T)-M(T')\ \ge\ m(\T)-m(\Sigma(T'))-\delta\bigl(2(n-1)\mu+2\mu_2\bigr),
\]
where $\mu_2$ is the expected number of pairs of splits proposed by a complete unit. Consequently the cutting-plane estimator is consistent for every $\delta<\gamma/\bigl(2(n-1)\mu+2\mu_2\bigr)$, a sufficient condition, where $\gamma>0$ is the complete-data margin of Theorem~\ref{thm:consistency}.
\end{enumerate}
\end{theorem}
\begin{proof}
\leanok
\uses{prop:buneman, thm:fragmentary, thm:consistency, thm:cutting}
(a) Let the members be non-trivial partial splits. If each has an extension and the extensions are pairwise compatible, the distinct extensions form a compatible set of full splits, which is the split set of a stemma (Proposition~\ref{prop:buneman}); resolving each vertex of degree above three by a binary subtree keeps every split, so a binary stemma displays each member. Conversely, if $T'$ displays every member, extending each member to a split of $T'$ that restricts to it gives pairwise compatible extensions, not necessarily distinct. The weight of a family is the sum of its members' weights, and adding a positive-weight partial split displayed by the same tree keeps the family completable, so a maximum-weight completable family consists of all positive-weight partial splits displayed by a maximiser of $W$. This is the display formulation of maximum compatibility for partial characters; it is used here, not claimed as new (see~(Steel 1992; Semple and Steel 2003)).

(b) $W(T')/L$ is an average of i.i.d.\ bounded variables with mean $M(T')$, and there are finitely many binary stemmata, so $W(T')/L\to M(T')$ for all of them simultaneously, almost surely. If $T$ is the unique maximum, eventually $W(T)>W(T')$ for every $T'\neq T$; if another binary $T'$ displayed the selected family, it would attain at least the optimal weight, a contradiction, so the family is displayed by $T$ alone. If $M(T')>M(T)$, then $W(T')>W(T)$ eventually, so $T$ is not selected. A unique maximum does not imply that every split of $T$ is ever observed on all witnesses: a split with $m_\sigma=0$ may still be displayed through the restrictions of other proposals. If every $m_\sigma>0$ and $\delta<1$, each split of $T$ is eventually observed in full.

(c) A partial proposal of a unit is $\sigma'|_{O_u}$ for some complete proposal $\sigma'$. For $T$: every true $\sigma$ proposed by a complete unit and non-trivial on $O_u$ is displayed, which happens with probability at least $(1-4\delta)m_\sigma$ (Theorem~\ref{thm:fragmentary}(b)); two true proposals of the same unit can merge into one restriction only if the witnesses that separate them are illegible, with probability at most $2\delta$. So $M(T)\ge(1-4\delta)m(\T)-2\delta\mu_2$. For $T'$: a proposal $\sigma'\in\Sigma(T')$ contributes at most $m_{\sigma'}$; a proposal $\sigma'\notin\Sigma(T')$ is displayed by $T'$ only if $\sigma'|_{O_u}=\tau|_{O_u}$ for some $\tau\in\Sigma(T')$, which needs all the witnesses of one of the two mismatch sets of $\sigma'$ and $\tau$ to be illegible, with probability at most $2\delta$ for each of the $n-3$ splits $\tau$. So $M(T')\le m(\Sigma(T'))+2(n-3)\delta\mu$. As in Theorem~\ref{thm:fragmentary}(b), the mismatch sets are non-empty, so each is illegible with probability at most $\delta$; and pendant edges of $T'$ cannot display a non-trivial restriction. Subtracting, and using $m(\T)\le\mu$, gives the inequality. Finally $m(\T)-m(\Sigma(T'))=m(A)-m(B)\ge\gamma|A|\ge\gamma$ with $A=\T\setminus\Sigma(T')$, $B=\Sigma(T')\setminus\T$, by the definition of $\gamma$, and (b) concludes.
\end{proof}

\begin{proposition}[A tolerance independent of the number of witnesses; Proposition 7.15 of the report]
\label{prop:dimfree}
\lean{Stemma.CuttingDimFree.four_parts, Stemma.CuttingDimFree.display_prob_le, Stemma.CuttingDimFree.display_prob_le_text, Stemma.CuttingDimFree.cutting_dimfree, Stemma.CuttingDimFree.cutting_dimfree_margin, Stemma.CuttingDimFree.cutting_text_margin}
\leanok
\uses{def:gamma, def:joint}
Under the illegibility model of Theorem~\ref{thm:fragmentary}, let $T'$ be a binary stemma and $\sigma'=A|B$ a non-trivial full split not in $\Sigma(T')$. The probability that $\sigma'|_{O_u}$ is non-trivial and displayed by $T'$ is at most $2\max(\delta,2\delta^2)$, hence at most $2\delta$ for $\delta\le\frac12$. Consequently, for $\delta\le\frac12$ and every binary $T'\ne T$,
\[
M(T)-M(T')\ \ge\ m(\T)-m(\Sigma(T'))-\delta\bigl(6\mu+2\mu_2\bigr),
\]
and the cutting-plane estimator is consistent for every $\delta<\gamma/(6\mu+2\mu_2)$, whatever the number of witnesses. Under the hypothesis of Proposition~\ref{prop:textillegible} the probability is at most $4\delta$, and the estimator is consistent for every $\delta<\gamma/(8\mu+2\mu_2)$.
\end{proposition}
\begin{proof}
\leanok
\uses{thm:cuttingconsistency}
For a legible set $O$, the restriction $\sigma'|_O=(A\cap O)|(B\cap O)$ is displayed by $T'$, with both sides of size at least two, if and only if $\mathrm{conv}_{T'}(A\cap O)$ and $\mathrm{conv}_{T'}(B\cap O)$ are disjoint: a split of $T'$ that restricts to it separates the two subtrees, and conversely any edge on the path between two disjoint subtrees has a split that restricts to it. As $\sigma'\notin\Sigma(T')$, the subtrees $\mathrm{conv}_{T'}(A)$ and $\mathrm{conv}_{T'}(B)$ meet; their intersection $I$ is a subtree whose vertices are internal vertices of $T'$, each with at least two of its three edges in each subtree, hence one edge in both, so $I$ has an edge and a leaf $u$. Let $e_3$ be the edge of $I$ at $u$ and $e_1,e_2$ the other edges of $T'$ at $u$, with leaf sets $X_1,X_2,X_3$ beyond them. Since $e_3$ lies in both subtrees, $X_3$ meets $A$ and $B$; since $e_1,e_2\notin I$, each of $X_1,X_2$ lies within $A$ or within $B$; and since $u$ has degree at least two in each subtree, one lies in $A$ and the other in $B$. So $A=A_1\sqcup A_2$ and $B=B_1\sqcup B_2$ with $A_1=X_1$, $B_1=X_2$, $A_2=X_3\cap A$, $B_2=X_3\cap B$, all non-empty. If $A_1\cap O$ and $A_2\cap O$ are non-empty, the path between a witness of each passes through $u$, so $u\in\mathrm{conv}_{T'}(A\cap O)$; likewise for $B$. Display therefore requires one of $A_1,A_2,B_1,B_2$ to be wholly illegible, while $|A\cap O|,|B\cap O|\ge2$. The event that $A_1$ is wholly illegible and $|A\cap O|\ge2$ is empty unless $|A_2|\ge2$, and has probability at most $\delta^{|A_1|}$; so the $A$ side contributes at most $\mathbf 1[|A_2|\ge2]\delta^{|A_1|}+\mathbf 1[|A_1|\ge2]\delta^{|A_2|}\le\max(\delta,2\delta^2)$, and the $B$ side the same. Under the hypothesis of Proposition~\ref{prop:textillegible}, each of the four events contains the illegibility of one fixed witness, whose conditional probability given the readings is at most $\delta$, which gives $4\delta$. In the proof of Theorem~\ref{thm:cuttingconsistency}(c), the bound $2(n-3)\delta\mu$ on the contribution of proposals outside $\Sigma(T')$ is thereby replaced by $2\delta\mu$, respectively $4\delta\mu$, and the rest of the proof is unchanged; since $\gamma\le\mu$, the condition on $\delta$ implies $\delta<\frac16$.
\end{proof}

\chapter{Non-binary stemmata and coincident error}\label{sec:bedier}

\section{The Bédier effect, a threshold and a rival test}

\begin{proposition}[Bédier effect; Proposition 8.1 of the report]
\label{prop:bediereffect}
\lean{Stemma.Polytomy.bedier_effect}
\leanok
\uses{def:stemma, def:collation}
Let $T$ be any stemma, and suppose some false split $\tau$ compatible with every split of $\T$ has $m_\tau>0$. Then $\Pr(\hat S_L=\T)\le(1-m_\tau)^L\to0$.
\end{proposition}
\begin{proof}
\leanok
If $w(\tau)>0$, the compatible set $\T\cup\{\tau\}$ outweighs $\T$, so $\T$ is not a solution of~\eqref{eq:mwis}. Hence $\hat S_L=\T$ requires that no unit propose $\tau$, an event of probability $(1-m_\tau)^L$.
\end{proof}

\begin{corollary}[Polytomies that survive; Corollary 8.2 of the report]
\label{cor:tops}
\lean{Stemma.RootedTop.two_branched_top, Stemma.RootedTop.rooted_bedier_true}
\leanok
\uses{def:rooted, def:collation}
In the rooted coding, let $\hat S$ be a solution of~\eqref{eq:mwis} for clusters, and let $V$, a cluster of $\hat S$ or $\X$ itself, have $k\ge3$ children $D_1,\dots,D_k$, the maximal members of $\hat S$ and singletons strictly inside $V$. Then no unit proposes $D_i\cup D_j$ for $i\neq j$. In particular, if the root of the true rooted stemma has three children $C_1,C_2,C_3$ and $\tau=C_i\cup C_j$ has $m_\tau>0$, the probability that the true clusters are a solution is at most $(1-m_\tau)^L$.
\end{corollary}
\begin{proof}
\leanok
\uses{prop:bediereffect, prop:rooted}
$D_i\cup D_j$ has at least two members, is not $\X$ since $k\ge3$, and is not in $\hat S$, since it would then be a member strictly inside $V$ containing $D_i$. Every cluster of $\hat S$ contains $V$, lies inside one child, or is disjoint from $V$, so it is laminar with $D_i\cup D_j$. If a unit proposed $D_i\cup D_j$, adding it to $\hat S$ would give a heavier laminar set. The second statement is Proposition~\ref{prop:bediereffect} for clusters: $\tau$ is laminar with every clade.
\end{proof}

\begin{theorem}[Bédier threshold; Theorem 8.3 of the report]
\label{thm:bedier}
\lean{Stemma.Chernoff.bedier_threshold}
\leanok
\uses{def:stemma, def:collation}
Let $T$ be any stemma with at least one non-trivial split, let units be i.i.d., and let $b=\min_{\sigma\in\T}m_\sigma>a=\max_{\tau\notin\T}m_\tau$. For $c\in(a,b)$ let $\hat S_{L,c}$ be the solution of~\eqref{eq:mwis} restricted to the splits with $w(\sigma)>cL$. Then
\[
\Pr(\hat S_{L,c}\neq\T)\le\sum_{\sigma\in\T}e^{-L\,\mathrm{KL}(c\,\|\,m_\sigma)}+\sum_{\tau\in\F}e^{-L\,\mathrm{KL}(c\,\|\,m_\tau)},
\]
where $\mathrm{KL}(x\|y)=x\log\frac xy+(1-x)\log\frac{1-x}{1-y}$. With $c=\bigl((\sqrt a+\sqrt b)/2\bigr)^2$ the bound is at most $(|\T|+K)\exp\{-L(\sqrt b-\sqrt a)^2/4\}$.
\end{theorem}
\begin{proof}
\leanok
\uses{prop:threshold}
If every true split has $w>cL$ and every false split $w\le cL$, only the true splits survive, they are compatible and have positive weight, and the restricted solution is $\T$ (Proposition~\ref{prop:threshold}). Each $w(\sigma)$ is binomial; the Chernoff bound gives $\Pr(w(\sigma)\le cL)\le e^{-L\mathrm{KL}(c\|m_\sigma)}$ for $m_\sigma>c$ and $\Pr(w(\tau)>cL)\le e^{-L\mathrm{KL}(c\|m_\tau)}$ for $m_\tau<c$~(Chernoff 1952); false splits with $m_\tau=0$ never occur. For the second form, $\mathrm{KL}(x\|y)\ge-2\log\bigl(\sqrt{xy}+\sqrt{(1-x)(1-y)}\bigr)\ge(\sqrt x-\sqrt y)^2$, and with the chosen $c$ each term is at least $(\sqrt b-\sqrt a)^2/4$.
\end{proof}

\begin{proposition}[Largest-gap threshold; Proposition 8.4 of the report]
\label{prop:gap}
\lean{Stemma.GapThreshold.gap_separation, Stemma.GapThreshold.gap_threshold}
\leanok
\uses{def:stemma, def:collation}
Sort the expected supports of $\T\cup\F$ together with $0$, and suppose the gap $G=b-a$ between the largest false and the smallest true value exceeds every other gap between consecutive values, the largest of which is $g_2$. Cut the observed values $w/L$, together with $0$, at the midpoint of their largest gap, or of any of them if several are equal. Then $\Pr(\text{the cut does not separate }\T\text{ from }\F)\le2(|\T|+K)\exp\{-L(G-g_2)^2/8\}$.
\end{proposition}
\begin{proof}
\leanok
If every $|w(\sigma)/L-m_\sigma|<\eta=(G-g_2)/4$, the observed values keep the order between the two groups, the observed gap between them exceeds $G-2\eta$, and every other observed gap is below $g_2+2\eta=G-2\eta$, while the gap between the groups exceeds it; splits with $m_\tau=0$ are observed at $0$. Hoeffding's inequality bounds each deviation by $2e^{-2L\eta^2}$. The deterministic step is checked on random instances (Appendix~B of the report).
\end{proof}

\begin{proposition}[Rival test; Proposition 8.5 of the report]
\label{prop:rival}
\lean{Stemma.Rival.rival}
\leanok
\uses{def:collation}
Let $\sigma$ and $\tau$ be in conflict, with i.i.d.\ units. Conditionally on $w(\sigma)+w(\tau)=N$, $w(\sigma)$ is binomial with parameters $N$ and $m_\sigma/(m_\sigma+m_\tau)$. Hence $\Pr\bigl(w(\sigma)\ge k\mid N\bigr)$ under $m_\sigma=m_\tau$ is an exact test of whether $\sigma$ is better supported than its rival.
\end{proposition}
\begin{proof}
\leanok
\uses{lem:exclusive}
By Lemma~\ref{lem:exclusive} no unit proposes both. Each unit independently proposes $\sigma$, $\tau$ or neither, with probabilities $m_\sigma$, $m_\tau$ and $1-m_\sigma-m_\tau$, so the counts are multinomial, and conditioning on their sum gives the binomial.
\end{proof}

\section{Coincident errors and the choice of coding}\label{sec:coincidence}

\begin{definition}[Attraction model]
\label{def:attraction}
\lean{Stemma.Mechanisms.read, Stemma.Mechanisms.PropU, Stemma.Mechanisms.PropR, Stemma.Attraction.wt, Stemma.Attraction.omega}
\leanok
\uses{def:rooted, def:stemma}
Root $T$ at a vertex (to root on an edge, subdivide it and give each half its own rate) and let each edge $e$, with the clade $C_e$ of witnesses below it, independently carry a potential innovation with probability $\pi_e$. Each innovation adopts an \emph{attractive} reading, the same for the whole unit, with probability $\alpha\in[0,1]$, independently of the other edges, and otherwise a reading of its own that occurs nowhere else. A witness reads the value of the lowest innovating edge above it, or the original reading if there is none. For $\alpha=1$ this is the irreversible model of~(Jurczak and Jurczak 2026), in which parallel innovations always coincide.
\end{definition}

\begin{theorem}[Mechanisms of false support; Theorem 8.7 of the report]
\label{thm:mechanisms}
\lean{Stemma.Mechanisms.retention_coincidence_report, Stemma.Mechanisms.overwrite_report}
\leanok
\uses{def:attraction}
In the attraction model on a rooted tree, let $M$ be the set of edges carrying an innovation in a unit.
\begin{enumerate}[(i)]
\item If $|M|\le1$, every proposed split or cluster is true or trivial, in both codings.
\item (\emph{Retention and coincidence.}) If $M=\{e,f\}$ with $C_e\cap C_f=\emptyset$, the reading classes are $C_e$, $C_f$ and the rest, with $C_e$ and $C_f$ merged if and only if both readings are the attractive one. The unrooted coding proposes $(C_e\cup C_f)|\text{rest}$ whatever the two readings, provided both sides have at least two witnesses, and the rooted coding proposes the cluster $C_e\cup C_f$ if and only if both readings are the attractive one and $C_e\cup C_f\ne\X$.
\item (\emph{Overwrite.}) If $M=\{e,f\}$ with $C_f\subsetneq C_e$, both codings propose $C_e\setminus C_f$ unless both readings are the attractive one, provided $|C_e\setminus C_f|\ge2$.
\end{enumerate}
In each case the other proposals are true or trivial. Consequently, for every split or cluster $\tau$,
\[
\Lambda_\tau\ \le\ m_\tau\ \le\ \Lambda_\tau+\Pr(|M|\ge3),\qquad \Pr(|M|\ge3)\le\tfrac16\Bigl(\sum_e\pi_e\Bigr)^3,
\]
where $\Lambda_\tau$, the contribution of at most two innovations, sums $\pi_e\prod_{g\ne e}(1-\pi_g)$ over single edges proposing $\tau$ and $\pi_e\pi_f\,\omega_{ef}(\tau)\prod_{g\ne e,f}(1-\pi_g)$ over pairs, with $\omega_{ef}(\tau)=\alpha^2\,\mathbf 1[\tau\text{ proposed when both readings are attractive}]+(1-\alpha^2)\,\mathbf 1[\tau\text{ proposed otherwise}]$. For a false target this is $1$ for retention in the unrooted coding, $\alpha^2$ for coincidence in the rooted coding and $1-\alpha^2$ for overwrite in both; for a true target, a disjoint pair $\{e,f\}$ proposes $\sigma_e$, and the cluster $C_e$, with weight $1-\alpha^2$, and a nested pair $C_f\subsetneq C_e$ proposes $\sigma_e$ with weight $1$ in the unrooted coding and the cluster $C_e$ with weight $\alpha^2$ in the rooted one, and proposes $\sigma_f$, or $C_f$, with weight $1-\alpha^2$.
\end{theorem}
\begin{proof}
\leanok
A witness reads the value of the lowest innovating edge above it. With one innovation on $e$, the classes are $C_e$ and its complement, which propose $\sigma_e$ or the cluster $C_e$. With two on disjoint clades, the classes are $C_e$, $C_f$ and the remaining witnesses, which keep the original reading, unless both readings are attractive, in which case $C_e\cup C_f$ is one class; the remaining witnesses propose $(C_e\cup C_f)|\text{rest}$ in the unrooted coding in either case, while in the rooted coding the original reading proposes nothing. With $C_f\subsetneq C_e$, the witnesses of $C_f$ take $f$'s reading and those of $C_e\setminus C_f$ keep $e$'s, so these are two classes unless the two readings coincide, which happens only when both are attractive. The bounds follow by summing over outcomes with at most two innovations, which are counted exactly, and bounding the rest by the probability of three or more; $\Pr(|M|\ge3)\le\sum_{|S|=3}\prod_{e\in S}\pi_e\le(\sum\pi_e)^3/6$. The classification is checked by machine on random trees (Appendix~B of the report).
\end{proof}

\begin{corollary}[Comparison of the codings; Corollary 8.8 of the report]
\label{cor:dominance}
\lean{Stemma.Dominance.dominance, Stemma.Dominance.omega_dominance}
\leanok
\uses{def:attraction}
Let $\tau=C|D$ be a non-trivial split. In every outcome of a unit, if the rooted coding proposes the cluster $C$ or the cluster $D$, the unrooted coding proposes $\tau$; hence, at every order, the support of each of $C$ and $D$ in the rooted coding is at most the support of $\tau$ in the unrooted coding. At second order, when $\tau$ is not a split of the tree, each disjoint pair proposing the cluster carries the factor $\alpha^2$ instead of $1$, nested pairs carry the same factor in both codings, and pairs whose union is the other side propose $\tau$ in the unrooted coding only.
\end{corollary}
\begin{proof}
\leanok
\uses{thm:mechanisms}
A cluster proposed in the rooted coding is the class of a reading of the unit, with at least two witnesses; the other side of $\tau$ has at least two witnesses because $\tau$ is non-trivial; so the unrooted coding proposes $\tau$ from the same reading. Taking expectations gives the inequality at every order. The second-order weights follow from Theorem~\ref{thm:mechanisms}: a cluster of a split that is not a split of the tree is proposed in the rooted coding only by a disjoint pair of innovations whose readings coincide, with weight $\alpha^2$, or by a nested pair, with weight $1-\alpha^2$; the same pairs propose $\tau$ in the unrooted coding with weights $1$ and $1-\alpha^2$, and a disjoint pair whose union is the other side proposes $\tau$ in the unrooted coding only.
\end{proof}

\begin{definition}[Irreversible model]
\label{def:irreversible}
\lean{Stemma.Irreversible.ones, Stemma.Irreversible.pM, Stemma.Irreversible.Proposes, Stemma.Irreversible.suppI}
\leanok
\uses{def:stemma, def:collation}
Root $T$ at a vertex (to root on an edge, subdivide it first and give each half its own rate) and let $\mathcal E$ be its edges, each $e$ with the clade $C_e\subseteq\X$ of witnesses below it; $C_e$ is \emph{non-trivial} if $2\le|C_e|\le n-2$. A unit has an original reading $0$ and one alternative $1$. Independently for each edge, a potential innovation is marked with probability $\pi_e$, and a witness reads $1$ if some edge above it is marked; no innovation is undone. The witnesses reading $1$ form $Y=\bigcup\{C_e:e\text{ marked}\}$, and the unit proposes the split $Y|\X\setminus Y$ if it is non-trivial.
\end{definition}

\begin{proposition}[Support bounds; Proposition 8.10 of the report]
\label{prop:modelbounds}
\lean{Stemma.Irreversible.suppI_edge_ge, Stemma.Irreversible.false_mass_sides}
\leanok
\uses{def:irreversible}
Under the irreversible model, for every edge $e$ with $C_e$ non-trivial,
\[
m_{\sigma_e}\ \ge\ \pi_e\prod_{f:\,C_f\not\subseteq C_e}(1-\pi_f),
\]
and the total expected support of false splits satisfies
\[
\sum_{\tau\notin\T}m_\tau\ \le\ \sum_{\{e,f\}:\,C_e\cap C_f=\emptyset}\pi_e\pi_f .
\]
\end{proposition}
\begin{proof}
\leanok
If $e$ is marked and no edge $f$ with $C_f\not\subseteq C_e$ is, then $Y=C_e$: marks on edges with $C_f\subseteq C_e$ add nothing. This event has exactly the stated probability, which gives the first bound. The marked edges whose clades are maximal among the marked ones have pairwise disjoint clades, by laminarity, and $Y$ is their union. If there are none or one, the proposed split is absent, trivial or true. A false split therefore requires two marked edges with disjoint clades, and since a unit proposes at most one split, the total false mass is at most the probability of that event, bounded by the sum over unordered pairs.
\end{proof}

\begin{corollary}[A sufficient condition; Corollary 8.11 of the report]
\label{cor:modelglobal}
\lean{Stemma.Irreversible.model_global, Stemma.Irreversible.model_consistency}
\leanok
\uses{def:irreversible, def:gamma}
If $T$ is binary and
\[
\min_{e:\ C_e\text{ non-trivial}}\pi_e\prod_{f\in\mathcal E}(1-\pi_f)\ >\ \sum_{C_e\cap C_f=\emptyset}\pi_e\pi_f,
\]
(the sum over unordered pairs), then global separation holds for $m$, hence $\gamma>0$ and the estimator is consistent.
\end{corollary}
\begin{proof}
\leanok
\uses{prop:modelbounds, thm:global, thm:consistency}
By Proposition~\ref{prop:modelbounds} the left side is at most every $m_{\sigma_e}$ (multiplying by the extra factors $1-\pi_f$ only lowers it) and the right side at least every false $m_\tau$. When the root subdivides an internal edge, its two halves give lower bounds for the same split.
\end{proof}

\begin{proposition}[Quartet; Proposition 8.12 of the report]
\label{prop:quartet}
\lean{Stemma.Quartet.m12, Stemma.Quartet.m13, Stemma.Quartet.m14}
\leanok
\uses{def:irreversible}
Let $n=4$, $T=12|34$, rooted on the internal edge with halves of rates $r_1$ (above $\{1,2\}$) and $r_2$ (above $\{3,4\}$), and pendant rates $p_1,\dots,p_4$. With $\bar x=1-x$,
\begin{align*}
m_{12|34}&=\bar r_2\bar p_3\bar p_4\,(r_1+\bar r_1p_1p_2)+\bar r_1\bar p_1\bar p_2\,(r_2+\bar r_2p_3p_4),\\
m_{13|24}&=\bar r_1\bar r_2\,(p_1p_3\,\bar p_2\bar p_4+p_2p_4\,\bar p_1\bar p_3),\\
m_{14|23}&=\bar r_1\bar r_2\,(p_1p_4\,\bar p_2\bar p_3+p_2p_3\,\bar p_1\bar p_4).
\end{align*}
The three splits are pairwise in conflict, so $\hat S_L$ is a most-supported split of positive support, or $\emptyset$ if none occurs. The estimator is consistent if and only if $m_{12|34}>\max(m_{13|24},m_{14|23})$.
\end{proposition}
\begin{proof}
\leanok
\uses{thm:consistency, lem:exclusive}
On four witnesses a unit proposes a non-trivial split exactly when $|Y|=2$. $Y=\{1,2\}$ requires no innovation on $r_2$, $p_3$, $p_4$, and either $r_1$, or not $r_1$ and both $p_1,p_2$; symmetrically for $\{3,4\}$. $Y=\{1,3\}$ requires $p_1,p_3$ and none of $r_1,r_2,p_2,p_4$; similarly for the other pairs. The three splits of a quartet cross pairwise. With $n-3=1$, $\gamma=m_{12|34}-\max(m_{13|24},m_{14|23},0)$ and Theorem~\ref{thm:consistency} gives consistency when $\gamma>0$ and inconsistency when a false split is strictly heavier. At a tie between $12|34$ and a false split with common positive value $a$, the per-unit difference takes the values $\pm1$ with probability $a$ each (Lemma~\ref{lem:exclusive}), so it has variance $2a>0$ and the false split strictly outscores the true one with probability tending to $\tfrac12$. If $m_{12|34}=0$, the true split is never selected.
\end{proof}

\begin{corollary}[Textual long-branch zone; Corollary 8.13 of the report]
\label{cor:lba}
\lean{Stemma.Quartet.long_branch, Stemma.CodingLimit.long_branch_consistent, Stemma.CodingLimit.long_branch_inconsistent}
\leanok
\uses{def:irreversible}
Let $p_2=p_4=0$ and $r_1=r_2=r<1$. The estimator is consistent if and only if
\[
(1-r)\,p_1p_3<r\,(2-p_1-p_3),
\]
and under the reverse strict inequality it selects the false split $13|24$ with probability tending to one.
\end{corollary}
\begin{proof}
\leanok
\uses{prop:quartet}
Then $m_{12|34}=r(1-r)(2-p_1-p_3)$, $m_{13|24}=(1-r)^2p_1p_3$ and $m_{14|23}=0$; divide by $1-r>0$ and apply Proposition~\ref{prop:quartet}. At $r=1$ every witness reads $1$ and nothing is ever recovered.
\end{proof}

\begin{theorem}[At leading order, coincidence enters only the rooted coding; Theorem 8.14 of the report]
\label{thm:coding}
\lean{Stemma.CodingLimit.unrooted_consistency, Stemma.CodingLimit.rooted_inconsistency, Stemma.CodingLimit.unrooted_false_selected, Stemma.CodingLimit.rooted_consistency, Stemma.CodingLimit.rooted_selects_13}
\leanok
\uses{def:attraction, def:gamma}
Let $n=4$ and $T=12|34$, rooted on the internal edge, whose two halves have equal rate $r$, with pendant rates $p_1,p_3$ and $p_2=p_4=0$. Then the false supports are, exactly,
\[
m^{\mathrm{unr}}_{13|24}=p_1p_3\bigl[(1-r)^2+2\alpha^2r(1-\alpha)(1-\alpha r)\bigr],\qquad
m^{\mathrm{root}}_{\{1,3\}}=\alpha^2p_1p_3(1-\alpha r)^2 .
\]
Let $r=\rho\varepsilon^2$, $p_1=a\varepsilon$, $p_3=b\varepsilon$ with $a,b,\rho>0$ fixed. For every sufficiently small fixed $\varepsilon>0$, as the number of units grows:
\begin{enumerate}[(a)]
\item in the unrooted coding, for every $\alpha$, the estimator is consistent if $ab<2\rho$ and selects the false split $13|24$ with probability tending to one if $ab>2\rho$;
\item in the rooted coding, it is consistent if $\alpha^2ab<2\rho$ and inconsistent if $\alpha^2ab>2\rho$.
\end{enumerate}
\end{theorem}
\begin{proof}
\leanok
\uses{thm:consistency, cor:lba}
The supports are computed by summing over the $3^6$ outcomes (none, attractive, own reading) of the six edges; the full polynomials, for all rates, are in \texttt{data/kappa\_quartet.json}, checked against the irreversible model at $\alpha=1$ and against simulation (Appendix~B of the report). In the unrooted coding, two witnesses with readings of their own leave the other two sharing the original reading, which proposes $13|24$ whatever $\alpha$; in the rooted coding, $\{1,3\}$ is proposed only when both innovations take the attractive reading and neither internal half-edge takes the attractive reading, which would put witness $2$ or $4$ in the class of $1$ and $3$, whence the factor $\alpha^2(1-\alpha r)^2$.

(a) Expanding in $\varepsilon$, $m^{\mathrm{unr}}_{12|34}=2\rho\varepsilon^2+O(\varepsilon^3)$, $m^{\mathrm{unr}}_{13|24}=ab\varepsilon^2+O(\varepsilon^3)$ and $m^{\mathrm{unr}}_{14|23}=O(\varepsilon^4)$. The three splits cross pairwise, so $\gamma=(2\rho-ab)\varepsilon^2+O(\varepsilon^3)$, positive or negative for small $\varepsilon$ according to the sign of $2\rho-ab$; Theorem~\ref{thm:consistency} concludes.

(b) The proper clusters on four witnesses are the six pairs and four triples. Expanding, $m_{\{1,2\}}$ and $m_{\{3,4\}}$ are $\rho\varepsilon^2+O(\varepsilon^3)$, $m_{\{1,3\}}=\alpha^2ab\varepsilon^2+O(\varepsilon^3)$, and every other cluster has support $O(\varepsilon^3)$ (a triple needs an internal innovation and a pendant one; the pairs $\{1,4\},\{2,3\},\{2,4\}$ need an internal innovation and another innovation). The true value is $2\rho\varepsilon^2+O(\varepsilon^3)$. A laminar family containing $\{1,3\}$ contains neither true cluster, which both conflict with it, so its value is $\alpha^2ab\varepsilon^2+O(\varepsilon^3)$; a laminar family containing neither $\{1,3\}$ nor both true clusters has value at most $\rho\varepsilon^2+O(\varepsilon^3)$. The margin is therefore positive for small $\varepsilon$ when $\alpha^2ab<2\rho$, and a family containing $\{1,3\}$ is strictly heavier when $\alpha^2ab>2\rho$. Theorem~\ref{thm:consistency}, which holds for clusters (Proposition~\ref{prop:rooted}), concludes.
\end{proof}

\begin{proposition}[Drift towards the root; Proposition 8.15 of the report]
\label{prop:drift}
\lean{Stemma.Drift.drift_unrooted_consistency, Stemma.Drift.drift_unrooted_false, Stemma.Drift.drift_rooted_consistency, Stemma.Drift.drift_rooted_false}
\leanok
\uses{def:attraction, def:gamma}
On the rooted tree $(((1,3),2),4)$, root at the vertex joining $\{1,2,3\}$ and $4$, let the edge above $\{1,3\}$ have rate $u\varepsilon^2$, the edge above $\{1,2,3\}$ rate $s\varepsilon$, the pendant edges of $3$ and $4$ rates $b\varepsilon$ and $t\varepsilon$, and those of $1$ and $2$ rate $0$, where $u,s>0$, $b,t\ge0$ and $\alpha\in[0,1]$ are fixed. For every sufficiently small fixed $\varepsilon>0$, as the number of units grows, the rooted coding is consistent if $(1-\alpha^2)sb<u$ and selects the false cluster $\{1,2\}$ if $(1-\alpha^2)sb>u$; the unrooted coding is consistent if $\bigl((1-\alpha^2)s+t\bigr)b<u$ and selects $12|34$ if the reverse strict inequality holds.
\end{proposition}
\begin{proof}
\leanok
\uses{thm:consistency, thm:mechanisms}
Exact enumeration (\texttt{code/drift\_quartet.py}, checked against simulation) gives, to leading order, $m_{13|24}=m_{\{1,3\}}=u\varepsilon^2$, $m_{12|34}=\bigl((1-\alpha^2)s+t\bigr)b\varepsilon^2$, $m_{\{1,2\}}=(1-\alpha^2)sb\varepsilon^2$, $m_{\{1,2,3\}}=s\varepsilon+O(\varepsilon^2)$, $m_{\{3,4\}}=\alpha^2tb\varepsilon^2$, and $O(\varepsilon^3)$ for every other split or cluster. In the unrooted coding the three splits cross pairwise and Theorem~\ref{thm:consistency} applies. In the rooted coding $\{1,2,3\}$, of order $\varepsilon$ since $s>0$, is in every optimum; it is compatible with $\{1,3\}$ and $\{1,2\}$, which conflict, and in conflict with $\{3,4\}$; the comparison is therefore between $m_{\{1,3\}}$ and $m_{\{1,2\}}$.
\end{proof}

\begin{proposition}[Two-reading coding; Proposition 8.16 of the report]
\label{prop:tworeading}
\lean{Stemma.TwoReading.two_reading_mechanisms, Stemma.TwoReading.two_reading_consistency, Stemma.TwoReading.two_reading_false_selected}
\leanok
\uses{def:attraction, def:gamma}
Let the \emph{two-reading coding} count a split at a unit only when its two sides are the two reading classes of the unit, so that a unit with three or more readings proposes nothing.
\begin{enumerate}[(i)]
\item In the attraction model on a rooted tree, a unit with at most one innovation proposes only true splits. With exactly two innovations, on $e$ and $f$, it proposes a split of no edge only when $C_e\cap C_f=\emptyset$ and both readings are the attractive one, and the split is then $(C_e\cup C_f)|\text{rest}$; a nested pair $C_f\subsetneq C_e$ proposes nothing, or the split of $e$ when both readings are attractive.
\item On the quartet of Theorem~\ref{thm:coding}, $m_{13|24}=\alpha^2p_1p_3(1-r)^2$, $m_{14|23}=0$ and $m_{12|34}=2r+O(\varepsilon^3)$, and with the scaling of the theorem the estimator is consistent if $\alpha^2ab<2\rho$ and selects $13|24$ with probability tending to one if $\alpha^2ab>2\rho$.
\end{enumerate}
\end{proposition}
\begin{proof}
\leanok
\uses{thm:mechanisms, thm:coding, thm:consistency}
(i) One innovation gives two classes, $C_e$ and the rest. Two innovations on disjoint clades give the classes $C_e$, $C_f$ and the rest, three unless both readings are attractive, when they are $C_e\cup C_f$ and the rest (if the rest is empty, $C_e|C_f$ is the split of the root's edges). A nested pair gives $C_f$, $C_e\setminus C_f$ and the rest, three unless both readings are attractive, when they are $C_e$ and the rest. (ii) With $p_2=p_4=0$, the classes are exactly $\{1,3\}$ and $\{2,4\}$ only when both halves are silent and both pendants take the attractive reading, with probability $\alpha^2(1-r)^2p_1p_3$; the class $\{1,4\}$ never arises; and a single innovation on one half proposes $12|34$, so $m_{12|34}=2r+O(\varepsilon^3)$. The three splits cross pairwise and Theorem~\ref{thm:consistency} concludes, as in the proof of Theorem~\ref{thm:coding}. The supports are checked by machine (Appendix~B of the report).
\end{proof}

\begin{proposition}[Several attractive readings; Proposition 8.17 of the report]
\label{prop:multistate}
\lean{Stemma.MultiAttraction.transfer_PU, Stemma.MultiAttraction.transfer_PR, Stemma.MultiAttraction.dominanceG}
\leanok
\uses{def:attraction}
Let each innovation take attractive reading $j\in\{1,\dots,k\}$ with probability $a_j\ge0$, independently of the other edges, or a reading of its own with probability $1-\sum_ja_j\ge0$, and let $A=\sum_ja_j^2$, the probability that two independent innovations take the same reading. With at most two innovations in a unit, the classes of witnesses, and the class that keeps the original reading, are those of the attraction model in which both innovations are attractive exactly when they take the same reading. Hence, for every split or cluster $\tau$, in either coding,
\[
\bigl|m_\tau-m^{(\sqrt A)}_\tau\bigr|\ \le\ \Pr(|M|\ge3)\ \le\ \tfrac16\Bigl(\sum_e\pi_e\Bigr)^3,
\]
where $m^{(\sqrt A)}$ is the expected support in the attraction model with $\alpha=\sqrt A$. Theorem~\ref{thm:mechanisms} and Propositions~\ref{prop:drift}, \ref{prop:alpha} and~\ref{prop:polytomynull} hold with $\alpha^2$ replaced by $A$, and so do parts (a) and (b) of Theorem~\ref{thm:coding}; Corollary~\ref{cor:dominance} holds unchanged.
\end{proposition}
\begin{proof}
\leanok
\uses{thm:mechanisms, thm:coding, prop:drift, prop:alpha, prop:polytomynull, cor:dominance}
With no innovation or one, the readings do not affect the classes. With two innovations on $e$ and $f$, every witness reads the value of the lowest innovation above it, so the classes are those of Theorem~\ref{thm:mechanisms}(ii)--(iii), and they depend on the two values only through whether they are equal. The values are equal exactly when both innovations take the same attractive reading, with probability $\sum_ja_j^2=A$, which is the probability $\alpha^2$ that both are attractive in the attraction model with $\alpha=\sqrt A$; since $A\le\sum_ja_j\le1$, that model is defined. The parts of $m_\tau$ and $m^{(\sqrt A)}_\tau$ with at most two innovations are therefore equal term by term, and each expected support exceeds that common part by at most $\Pr(|M|\ge3)$, which depends on the rates alone. The listed statements use the model only through these two-innovation parts and this bound, with remainders of order $\varepsilon^3$; Corollary~\ref{cor:dominance} holds outcome by outcome, whatever the readings.
\end{proof}

\section{Estimating the coincidence probability}\label{sec:alpha}

\begin{definition}[Isolated pair]
\label{def:isolated}
\lean{Stemma.Alpha.Isolated}
\leanok
\uses{def:attraction}
Call a pair of leaves $i,j$ \emph{isolated} in a rooted tree if $\{i,j\}$ is not a clade, no clade minus a sub-clade equals $\{i,j\}$ or its complement $X\setminus\{i,j\}$, no union of two disjoint clades other than $\{i\}\cup\{j\}$ equals either, and $X\setminus\{i,j\}$ is not a clade. Let $R_{ij}$ be the probability that a unit proposes the split $\{i,j\}|\text{rest}$ in the unrooted coding, and $C_{ij}$ the probability that it proposes the cluster $\{i,j\}$ in the rooted coding, with the original reading as reference.
\end{definition}

\begin{proposition}[Estimating $\alpha$; Proposition 8.18 of the report]
\label{prop:alpha}
\lean{Stemma.Alpha.ratio_tendsto, Stemma.Alpha.alpha_pooled}
\leanok
\uses{def:isolated}
In the attraction model on a rooted tree, for every isolated pair, $|C_{ij}-\alpha^2R_{ij}|\le\Pr(|M|\ge3)$. Hence, over a set $\mathcal P$ of isolated pairs with $\sum_{\mathcal P}R_{ij}>0$, the ratio of the observed counts $\sum\hat C_{ij}/\sum\hat R_{ij}$ converges to $\sum C_{ij}/\sum R_{ij}$, which differs from $\alpha^2$ by at most $|\mathcal P|\Pr(|M|\ge3)/\sum_{\mathcal P}R_{ij}$; with $n$ fixed, every rate at most of order $\varepsilon$ and the two pendant rates of every pair in $\mathcal P$ at least $c\varepsilon$ for a fixed $c>0$, the denominator is of order $\varepsilon^2$ and the bound is $O(\varepsilon)$. The ratio is unchanged if units that propose nothing are removed, so it does not depend on which invariant places a collation records.
\end{proposition}
\begin{proof}
\leanok
\uses{thm:mechanisms}
By Theorem~\ref{thm:mechanisms}, a single innovation proposes only a clade and its complement, neither of which is $\{i,j\}$ or its complement by isolation; and among pairs of innovations, isolation leaves only the pendant edges of $i$ and $j$. That pair proposes $\{i,j\}|\text{rest}$ in the unrooted coding whatever the two readings, and the cluster $\{i,j\}$ in the rooted coding exactly when both are attractive. The two-innovation parts are therefore $R^{(2)}_{ij}$ and $C^{(2)}_{ij}=\alpha^2R^{(2)}_{ij}$, and each probability exceeds its two-innovation part by at most $\Pr(|M|\ge3)$. The identity $C^{(2)}=\alpha^2R^{(2)}$ is checked exactly on random trees (Appendix~B of the report).
\end{proof}

\section{A calibrated Bédier test}\label{sec:calibrated}

\begin{proposition}[Coincidence null at a polytomy; Proposition 8.19 of the report]
\label{prop:polytomynull}
\lean{Stemma.PolytomyNull.null_bound, Stemma.PolytomyMix.null_mixture}
\leanok
\uses{def:attraction}
In the attraction model on a rooted tree, with every innovation probability of order $\varepsilon$, let a vertex $v$ have children whose clades are $C_1,\dots,C_k$, $k\ge3$, and let $\tau=C_i\cup C_j$ with $i\ne j$. For a set $D$ of witnesses, singletons included, let $q_D$ be the probability that $D$ is exactly the set of witnesses carrying some non-archetype reading. In the rooted coding,
\[
m_\tau=\alpha^2\,q_{C_i}\,q_{C_j}+\mathbf 1[k=3,\ v\ \text{not the root}]\,(1-\alpha^2)\,q_{C_v}\,q_{C_l}+O(\varepsilon^3),
\]
where $C_l$ is the third clade when $k=3$. Suppose instead that, given a factor $\xi\in[0,\xi_{\max}]$ drawn independently for each unit, the edges innovate independently with probabilities $\xi\pi_e\le1$; then the two products on the right are each multiplied by $\kappa=\mathbb E\xi^2/(\mathbb E\xi)^2$, and the remainder is $O(\varepsilon^3\mathbb E\xi^3)$.
\end{proposition}
\begin{proof}
\leanok
\uses{thm:mechanisms}
The cluster $\tau$ is not a clade, so by Theorem~\ref{thm:mechanisms}(i) no single innovation proposes it. Two innovations on disjoint clades $C_e,C_f$ propose $C_e\cup C_f$ in the rooted coding when both readings are attractive; every clade inside $\tau$ lies inside $C_i$ or $C_j$, so $C_e\cup C_f=\tau$ forces $\{e,f\}=\{e_i,e_j\}$, the edges above $C_i$ and $C_j$. Two nested innovations $C_f\subsetneq C_e$ propose $C_e\setminus C_f$ unless both readings are attractive. Here $C_e\supsetneq\tau$ is a clade, so $C_e\supseteq C_v$. If $C_e\supsetneq C_v$, then $C_e\setminus\tau$ contains a child clade $C_l$, $l\ne i,j$, and a witness outside $C_v$, and a clade containing both contains $C_v$, hence $\tau$; so $C_e\setminus\tau$ is not a clade. If $C_e=C_v$, then $C_v\setminus\tau$ is the union of the $k-2$ remaining child clades, which is a clade only when $k=3$; and when $v$ is the root there is no edge above it. The two-innovation part of $m_\tau$ is therefore $\alpha^2\pi_{e_i}\pi_{e_j}+\mathbf 1[\cdot](1-\alpha^2)\pi_{e_v}\pi_{e_l}$, up to a factor $1+O(\varepsilon)$. A single innovation on $e$ makes $C_e$ exactly a non-archetype class, and every other outcome that does so has two or more innovations, so $q_{C_e}=\pi_e+O(\varepsilon^2)$ for every clade, singletons included; the products of these first-order terms agree with the products of the $\pi$ to $O(\varepsilon^3)$. Given $\xi$ the two-innovation parts are $\xi^2$ times the products of the $\pi$; averaging, they become $\mathbb E\xi^2$ times those products, while $q_{C_e}=\mathbb E\xi\,\pi_e+O(\varepsilon^2)$, so each product of $q$'s must be multiplied by $\kappa$; outcomes with three or more innovations have probability $O(\varepsilon^3\mathbb E\xi^3)$ by a union bound over triples.
\end{proof}

\chapter{The graphs of stemmatics}\label{sec:graphs}

\section{What is forced}

\begin{proposition}[Size; Proposition 9.1 of the report]
\label{prop:size}
\lean{Stemma.Size.card_nontrivial}
\leanok
\uses{def:split}
On $n$ witnesses there are $2^{n-1}-n-1$ non-trivial full splits, and every independent set of the full conflict graph has at most $n-3$ vertices. Hence a best-supported stemma has at most $n-3$ families, however many candidates the variants propose. With fragmentary witnesses, candidates are indexed by sides and domain, so a pairwise relaxation can select many more partial splits; a jointly completable family still completes to at most $n-3$ distinct full splits.
\end{proposition}
\begin{proof}
\leanok
\uses{lem:facts}
There are $2^{n-1}-1$ unordered bipartitions of $\X$ into two non-empty sides, and $n$ of them have a side with one witness, which leaves $2^{n-1}-n-1$. The bound on independent sets is Lemma~\ref{lem:facts}(i); the completions of a jointly completable family are pairwise compatible full splits, and the same bound applies to the distinct ones.
\end{proof}

\begin{proposition}[Crossing set; Proposition 9.2 of the report]
\label{prop:crossing}
\lean{Stemma.Crossing.crossing_iff}
\leanok
\uses{def:stemma}
Let $T$ be any stemma and $\tau=A|B$ a non-trivial split. For an edge $e$ of $T$, $\tau$ is in conflict with the split of $e$ if and only if $e$ lies in both $T[A]$ and $T[B]$, the smallest subtrees spanning $A$ and $B$. Hence the edges whose splits conflict with $\tau$ form a connected subtree of internal edges; if $T$ is binary and $\tau\notin\Sigma(T)$, it is non-empty.
\end{proposition}
\begin{proof}
\leanok
\uses{lem:facts}
Let $e$ separate the leaves into $P|Q$. The split of $e$ conflicts with $\tau$ when $A$ and $B$ both meet both $P$ and $Q$, that is, when $e$ separates two witnesses of $A$ and two of $B$; an edge separates two witnesses of $A$ exactly when it lies on the path between them, that is, in $T[A]$. The intersection of two subtrees of a tree is connected, and an edge whose endpoints lie in both subtrees lies in both, since each subtree contains the path between its vertices. A pendant edge has a single witness on one side, so it cannot separate two witnesses of $A$ and two of $B$. Non-emptiness for binary $T$ is Lemma~\ref{lem:facts}(ii).
\end{proof}

\begin{proposition}[Five witnesses; Proposition 9.3 of the report]
\label{prop:five}
\lean{Stemma.Five.conflict_not_chordal, Stemma.Five.conflict_not_perfect}
\leanok
\uses{def:split}
On five witnesses the conflict graph of the ten non-trivial splits is isomorphic to the line graph of $K_5$, the complement of the Petersen graph. For every $n\ge5$, the conflict graph of all non-trivial splits contains it as an induced subgraph, and so is neither chordal nor perfect.
\end{proposition}
\begin{proof}
\leanok
For $n=5$, each non-trivial split is $ab|cde$, identified with the pair $\{a,b\}$. Two such splits conflict exactly when their pairs share one witness: if $\{a,b\}\cap\{c,d\}=\emptyset$ the sides $\{a,b\}$ and $\{c,d\}$ are disjoint; if they share $a$, the four intersections contain $a$, $b$, $c$ and the fifth witness. This is the line graph of $K_5$. For $n\ge5$, the splits $\{a,b\}|\X\setminus\{a,b\}$ with $a,b$ among five fixed witnesses conflict by the same rule. The line graph of $K_5$ has the induced $4$-cycle $12,23,34,41$ and the induced $5$-cycle $12,23,34,45,51$, so it is neither chordal nor perfect.
\end{proof}

\begin{theorem}[Fixed-parameter tractability in $n$; Theorem 9.4 of the report]
\label{thm:fpt}
\lean{Stemma.Dynamic.dp_correct}
\leanok
\uses{def:collation}
Fix a witness $x$. With the candidate weights, in $\mathbb N$, stored in a table indexed by the subsets of $\X\setminus\{x\}$ (zero off the candidates), the maximum weight of a compatible set of non-trivial full splits can be computed in $O(3^{n})$ additions and comparisons of weights, whatever the number of candidates. Building the table, enumerating subsets and arithmetic on long weights are not counted.
\end{theorem}
\begin{proof}
\leanok
Identify each split with its side avoiding $x$, a subset of $Y=\X\setminus\{x\}$ of size between $2$ and $n-2$. Two splits are compatible exactly when these sides are nested or disjoint, since the side containing $x$ of each meets the other's $x$-side. So compatible sets are laminar families on $Y$. For $C\subseteq Y$ let $F(C)$ be the maximum weight of a laminar family of subsets of $C$ that may include $C$, and $P(C)$ the maximum of $\sum_iF(D_i)$ over partitions of $C$ into blocks $D_i$. The maximal members of a laminar family inside $C$ other than $C$ are disjoint, so $F(C)=\max(0,w(C))+\max_{D}\bigl[F(D)+P(C\setminus D)\bigr]$, the maximum over $D\subsetneq C$ containing the least element of $C$, and $P(C)=\max_{D\ni\min C,\,D\subseteq C}[F(D)+P(C\setminus D)]$, with the conventions $w(C)=0$ unless $C$ is a candidate side ($2\le|C|\le n-2$, so in particular $w(Y)=0$), $F(\emptyset)=P(\emptyset)=0$, $F(\{i\})=0$, and an empty maximum equal to $0$; the partitions are thereby well defined down to singletons. Each recursion ranges over pairs $D\subseteq C\subseteq Y$, $3^{n-1}$ in all. The answer is $F(Y)$.
\end{proof}
